\section{Conclusion}

This report describes (read, will describe) a library for automatically deriving succinct data structures using the Haskell type system. By separating structural topology from data payloads and encoding the topology using DFUDS (Depth-First Unary Degree Sequence), we obtain a compact representation of trees that avoids traditional pointer-heavy layouts. This separation allows traversals to operate over contiguous memory rather than following chains of heap pointers. As a result, many operations reduce to cache-friendly operation over bitvectors, dramatically improving memory locality and reducing garbage collection overhead.

\subsection{Current Limitations}
The current prototype has several limitations that will be addressed in the next development phase:
\begin{itemize}
  \item \textbf{Only in the beta}, the DFUDS construction is implemented in a sub-optimal way, which is asymptotically inefficient. Our final version will construct the topology bitvector directly, without going through intermediate lists or repeated vector updates.
  \item The implementation currently supports only explicitly defined binary trees rather than arbitrary algebraic tree structures derived automatically via metaprogramming or any alternative.

  \item Navigational operations over the succinct representation remain limited. In particular, we aim to implement a Succinct Tree Zipper.
\end{itemize}

\subsection{Future Work}
Several directions remain for extending this work:
\begin{itemize}
  \item \textbf{Generic Derivation}: Extend the system to automatically support arbitrary algebraic data types by deriving succinct representations from sums and products.

  \item \textbf{Computational Logic Applications}: Evaluate the library in practical settings such as model checking and measuring whether succinct representations improve memory usage and processing bounds.

  \item \textbf{LOUDS Representation (if we have enough time)}: Introduce support for the Level Order Unary Degree Sequence (LOUDS) representation to enable efficient breadth-first navigation and accelerate BFS-style algorithms.
\end{itemize}